% Mathématiques 5e — cours V3
% Puissances : carré et cube

\section{Objectifs}

\begin{itemize}
\item Découvrir la notion de puissance.
\item Comprendre les notations carré et cube.
\item Connaître les carrés des entiers de $0$ à $12$.
\item Connaître le cube de $10$.
\item Écrire certains nombres sous la forme d'un carré ou d'un cube.
\item Calculer des expressions numériques contenant des puissances simples.
\item Calculer la valeur d'une expression littérale contenant une puissance simple.
\item Relier carré, cube, aire et volume.
\end{itemize}

\section{1. Une multiplication répétée}

Une puissance permet d'écrire plus simplement une multiplication d'un nombre par lui-même.

Par exemple :
\[
5\times5
\]
se note :
\[
5^2.
\]

Et :
\[
3\times3\times3
\]
se note :
\[
3^3.
\]

En 5e, on travaille principalement avec les exposants :
\[
2
\]
et :
\[
3.
\]

\section{2. Le carré d'un nombre}

\begin{definition}
Le carré d'un nombre $a$ est le produit de ce nombre par lui-même :
\[
\boxed{a^2=a\times a}.
\]
\end{definition}

Exemples :
\[
4^2=4\times4=16,
\]
\[
7^2=7\times7=49,
\]
\[
0{,}5^2=0{,}5\times0{,}5=0{,}25.
\]

Le petit nombre $2$ placé en haut à droite est l'exposant.

\section{3. Le cube d'un nombre}

\begin{definition}
Le cube d'un nombre $a$ est le produit de trois facteurs égaux à $a$ :
\[
\boxed{a^3=a\times a\times a}.
\]
\end{definition}

Exemples :
\[
2^3=2\times2\times2=8,
\]
\[
4^3=4\times4\times4=64.
\]

\section{4. Pourquoi les mots carré et cube ?}

Les noms viennent de la géométrie.

\coursefigure{/images/mathematiques/5eme/puissances-carre-cube.svg}{Un carré de côté a et un cube de côté a illustrant les notations a au carré et a au cube.}{Le carré intervient naturellement dans une aire ; le cube intervient naturellement dans un volume.}

Pour un carré de côté $a$, l'aire vaut :
\[
a\times a=a^2.
\]

Pour un cube d'arête $a$, le volume vaut :
\[
a\times a\times a=a^3.
\]

\section{5. Lire correctement une puissance}

L'écriture :
\[
6^2
\]
se lit :
\[
\text{« six au carré »}
\]
ou :
\[
\text{« six puissance deux »}.
\]

L'écriture :
\[
6^3
\]
se lit :
\[
\text{« six au cube »}
\]
ou :
\[
\text{« six puissance trois »}.
\]

\begin{attention}
\[
6^2\neq6\times2.
\]

En effet :
\[
6^2=36
\]
alors que :
\[
6\times2=12.
\]
\end{attention}

\section{6. Les carrés de zéro à douze}

Les carrés suivants doivent devenir des automatismes.

\begin{tabular}{c|c}
Nombre & Carré \\
$0$ & $0$ \\
$1$ & $1$ \\
$2$ & $4$ \\
$3$ & $9$ \\
$4$ & $16$ \\
$5$ & $25$ \\
$6$ & $36$ \\
$7$ & $49$ \\
$8$ & $64$ \\
$9$ & $81$ \\
$10$ & $100$ \\
$11$ & $121$ \\
$12$ & $144$ \\
\end{tabular}

\section{7. Construire les carrés mentalement}

On peut retrouver certains carrés en s'appuyant sur des calculs connus.

Par exemple :
\[
11^2=11\times11.
\]

On peut écrire :
\[
11\times11=11\times(10+1).
\]

Donc :
\[
110+11=121.
\]

Ainsi :
\[
\boxed{11^2=121}.
\]

\section{8. Le cube de dix}

Le programme demande de connaître :
\[
10^3.
\]

On calcule :
\[
10^3=10\times10\times10=1000.
\]

Donc :
\[
\boxed{10^3=1000}.
\]

Cette égalité relie directement la puissance au volume d'un cube de côté $10$.

\section{9. Reconnaître un carré}

Certains nombres peuvent être écrits comme carrés d'entiers.

Par exemple :
\[
64=8^2.
\]

De même :
\[
81=9^2.
\]

Et :
\[
144=12^2.
\]

Reconnaître ces écritures permet de gagner du temps dans de nombreux calculs.

\section{10. Reconnaître un cube simple}

Exemples :
\[
8=2^3,
\]
\[
27=3^3,
\]
\[
64=4^3,
\]
\[
125=5^3.
\]

Même si tous ces cubes ne sont pas à mémoriser comme automatismes officiels, savoir reconstruire leur valeur aide à comprendre la notation.

\section{11. Priorités opératoires et puissances}

Lorsqu'une expression contient une puissance, on calcule la puissance avant les multiplications et additions qui l'entourent.

\coursefigure{/images/mathematiques/5eme/puissances-expression.svg}{Expression numérique comportant un carré, une multiplication et une addition.}{La puissance est évaluée avant le produit puis avant la somme finale.}

Considérons :
\[
A=3^2+4\times5.
\]

On calcule d'abord :
\[
3^2=9.
\]

Puis :
\[
4\times5=20.
\]

Enfin :
\[
A=9+20=29.
\]

Donc :
\[
\boxed{A=29}.
\]

\section{12. Exemple avec parenthèses}

Calculons :
\[
B=(2+3)^2.
\]

On calcule d'abord la parenthèse :
\[
2+3=5.
\]

Puis le carré :
\[
5^2=25.
\]

Ainsi :
\[
\boxed{B=25}.
\]

\begin{attention}
\[
(2+3)^2
\]
n'est pas égal à :
\[
2+3^2.
\]
\end{attention}

\section{13. Exemple avec produit}

Calculons :
\[
C=2\times5^2.
\]

On commence par :
\[
5^2=25.
\]

Puis :
\[
C=2\times25=50.
\]

Donc :
\[
\boxed{C=50}.
\]

\section{14. Exemple avec soustraction}

Calculons :
\[
D=10^2-6^2.
\]

On obtient :
\[
10^2=100,
\]
et :
\[
6^2=36.
\]

Donc :
\[
D=100-36=64.
\]

Ainsi :
\[
\boxed{D=64}.
\]

\section{15. Valeur d'une expression littérale}

Une expression littérale peut contenir une puissance.

Par exemple :
\[
E=x^2+3.
\]

Pour :
\[
x=4,
\]
on remplace $x$ par $4$ :
\[
E=4^2+3.
\]

On calcule :
\[
4^2=16,
\]
puis :
\[
16+3=19.
\]

Donc :
\[
\boxed{E=19}.
\]

\section{16. Importance des parenthèses lors d'une substitution}

Si :
\[
F=x^2+2x,
\]
et :
\[
x=3,
\]
on écrit :
\[
F=3^2+2\times3.
\]

On obtient :
\[
F=9+6=15.
\]

Lorsque la valeur substituée comporte un signe ou une expression, les parenthèses deviennent essentielles.

\section{17. Aire d'un carré}

Un carré de côté :
\[
7\ \text{cm}
\]
a pour aire :
\[
7^2=49.
\]

Donc :
\[
\boxed{49\ \text{cm}^2}.
\]

L'unité d'aire est elle-même au carré :
\[
\text{cm}^2.
\]

\section{18. Volume d'un cube}

Un cube d'arête :
\[
4\ \text{cm}
\]
a pour volume :
\[
4^3=64.
\]

Donc :
\[
\boxed{64\ \text{cm}^3}.
\]

L'unité de volume est au cube :
\[
\text{cm}^3.
\]

\section{19. Comparer carré et double}

Le carré et le double sont deux opérations différentes.

Pour :
\[
a=6,
\]
le double vaut :
\[
2a=12.
\]

Le carré vaut :
\[
a^2=36.
\]

Il faut donc identifier précisément l'expression demandée.

\section{20. Programme de calcul}

Considérons le programme :
\begin{enumerate}
\item choisir un nombre ;
\item calculer son carré ;
\item ajouter $5$.
\end{enumerate}

Avec le nombre :
\[
4,
\]
on obtient :
\[
4^2+5.
\]

Puis :
\[
16+5=21.
\]

La sortie vaut :
\[
\boxed{21}.
\]

\section{21. Méthode de calcul}

\begin{methode}
Pour calculer une expression comportant des carrés ou des cubes :
\begin{enumerate}
\item traiter les parenthèses si nécessaire ;
\item calculer les puissances ;
\item effectuer multiplications et divisions ;
\item terminer par additions et soustractions ;
\item contrôler l'ordre de grandeur.
\end{enumerate}
\end{methode}

\section{22. Méthode de substitution}

\begin{methode}
Pour calculer une expression littérale contenant une puissance :
\begin{enumerate}
\item recopier l'expression ;
\item remplacer la lettre par sa valeur ;
\item ajouter des parenthèses si nécessaire ;
\item calculer les puissances ;
\item poursuivre selon les priorités opératoires.
\end{enumerate}
\end{methode}

\section{23. Erreurs fréquentes}

\begin{itemize}
\item Confondre $a^2$ avec $2a$.
\item Confondre $a^3$ avec $3a$.
\item Oublier de calculer la puissance avant une addition.
\item Croire que l'exposant est un facteur ordinaire.
\item Oublier les unités carrées pour une aire.
\item Oublier les unités cubiques pour un volume.
\item Remplacer une lettre sans respecter les parenthèses.
\end{itemize}

\section{À retenir}

\begin{aretenir}
Le carré d'un nombre est :
\[
\boxed{a^2=a\times a}.
\]

Le cube d'un nombre est :
\[
\boxed{a^3=a\times a\times a}.
\]

Les carrés des entiers de $0$ à $12$ sont à connaître.

Il faut également connaître :
\[
\boxed{10^3=1000}.
\]

Dans un calcul, les puissances sont évaluées avant les multiplications et les additions.

En géométrie :
\[
\boxed{\text{aire du carré}=a^2}
\]
et :
\[
\boxed{\text{volume du cube}=a^3}.
\]
\end{aretenir}
